% AI & Agent Tooling Stack — project report (generated documentation)
\documentclass[11pt,a4paper]{article}
\usepackage[margin=2.4cm]{geometry}
\usepackage[T1]{fontenc}
\usepackage{xcolor}
\usepackage{enumitem}
\usepackage{titlesec}
\usepackage[hidelinks]{hyperref}
\usepackage{listings}
\usepackage{graphicx}
\usepackage{float} % Added to force exact image placement

\definecolor{noc}{HTML}{0F766E}
\definecolor{nocdark}{HTML}{134E4A}

\titleformat{\section}{\large\bfseries\color{nocdark}}{\thesection}{0.8em}{}
\titleformat{\subsection}{\normalsize\bfseries\color{noc}}{\thesubsection}{0.8em}{}

\lstset{
  basicstyle=\ttfamily\small,
  breaklines=true,
  frame=single,
  rulecolor=\color{noc},
  backgroundcolor=\color[HTML]{F0FDFA},
}

\begin{document}

\begin{center}
  {\LARGE\bfseries AI \& Agent Tooling Stack}\\[4pt]
  {\large Local Model Hosting, LLM Routing, Browser Automation, Multi-Agent Orchestration}\\[8pt]
  {\color{noc}\rule{\textwidth}{1pt}}\\[6pt]
  {\small Marouane Aabirrouche --- IT Infrastructure specialist --- R\'eseaux \& Syst\`emes}\\
  {\small Project write-up PRJ-07 \,\textbullet\, work in progress --- living documentation}
\end{center}

\section*{Status note}
This node groups four small, individually undocumented projects into one
infrastructure-and-orchestration story: \emph{the tooling layer behind how I
build}. It is deliberately presented as in-progress rather than a finished,
shipped product.

\section{Overview}
Rather than treating AI tooling as a set of demos, this stack treats it as an
infrastructure problem: models are workloads with hardware constraints, requests
need a routing layer with policy, browsers become automatable endpoints, and
complex tasks decompose into agent teams with explicit delegation rules.

\section{Local AI hosting and quantization}
Consumer hardware sets real constraints: an 8\,GB VRAM budget (RTX~2060 Super)
with 32\,GB of system RAM. The work here is making useful models fit that box:

\begin{figure}[H]
  \centering
  \includegraphics[width=\textwidth]{unsloth.png}
  \caption{Unsloth interface hosting local quantized models.}
\end{figure}

\begin{itemize}[leftmargin=1.6em]
  \item Running \textbf{quantized} and smaller model variants so weights fit in
        VRAM while activations spill gracefully to system RAM.
  \item Evaluating \textbf{Mixture-of-Experts} architectures, whose sparse
        activation makes them unusually friendly to tight memory budgets ---
        explored via \textbf{llama.cpp}, \textbf{ExLlamaV2}/TabbyAPI, and
        \textbf{Ollama}, each with different trade-offs in speed, memory layout
        and serving ergonomics.
\end{itemize}

\section{OmniRoute --- self-hosted LLM routing layer}
OmniRoute is the piece that turns local inference into infrastructure. It runs
locally on \texttt{localhost:20128} and exposes an \textbf{OpenAI-compatible
API}, so any client that can talk to OpenAI can talk to the whole local fleet
without knowing what is behind it.

\begin{figure}[H]
  \centering
  \includegraphics[width=\textwidth]{home_omni.png}
  \caption{OmniRoute Provider Topology mapping connected backends.}
\end{figure}

Requests are routed across named model pools, each tuned for a job family:

\begin{lstlisting}
planning   coding   review/security   vision   cheap
\end{lstlisting}

\begin{figure}[H]
  \centering
  \includegraphics[width=\textwidth]{omni_combo.png}
  \caption{OmniRoute Combos establishing priority routing for different tasks.}
\end{figure}

The routing layer keeps callers decoupled from specific models: a pool can gain,
lose or swap backends without touching client code --- the same abstraction
discipline used for services behind a load balancer in classical infrastructure.

\section{Browser automation}
Web-based workflows were evaluated and operationalized with browser-driving
tools --- \textbf{Playwright MCP} and \textbf{Browser Use} --- giving agents
reliable hands on real pages: navigation, extraction, form flows and end-to-end
verification steps inside AI-assisted development loops.

\begin{figure}[H]
  \centering
  \includegraphics[width=\textwidth]{playwrit.png}
  \caption{MCP server configuration file integrating Playwright into the local environment.}
\end{figure}

\section{CrewAI multi-agent framework}
The orchestration layer is a hierarchical multi-agent system built on
\textbf{CrewAI}: \textbf{ten specialist agents} under a \textbf{custom manager
agent} with strict delegation-only behavior. The manager never does specialist
work itself; it decomposes objectives, assigns them, and integrates results.

\begin{figure}[H]
  \centering
  \includegraphics[width=\textwidth]{crew_ai.jpg}
  \caption{CrewAI multi-agent topology showing the Manager Agent delegating tasks to 10 specialists.}
\end{figure}

Model access flows through OmniRoute, so each role draws on the pool that fits
its task --- heavy reasoning for planning, cheap fast models for routine steps.

\section{Why it matters}
\begin{itemize}[leftmargin=1.6em]
  \item Treats AI capability as capacity planning under hardware constraints,
        not magic.
  \item A routing layer with named pools is the difference between a pile of
        models and an operable platform.
  \item Delegation-only managers encode accountability into agent systems ---
        clear ownership of every subtask.
\end{itemize}

\section{Next steps}
Documentation is intentionally thin today; planned follow-ups include publishing
configuration notes for OmniRoute's pool definitions, quantization benchmarks
against the 8\,GB VRAM budget, and a write-up of the CrewAI delegation patterns.

\vfill
\begin{center}
  {\small\color{noc}\rule{0.5\textwidth}{0.6pt}}\\[4pt]
  {\footnotesize Generated documentation --- NOC portfolio, PRJ-07. In progress; details evolve.}\\
  {\footnotesize Bouznika, Morocco \,\textbullet\, \today}
\end{center}

\end{document}